% Section 3.5, from lectures/L03/appendix/d_exercises.md.

\section{Exercises}
\label{c3:sec:exercises}

\begin{aside}
\textbf{How to check your work.} Exercises \ref{c3:ex:driver} and \ref{c3:ex:handler} are checked
by this chapter's test suite: write the file at the path the specification gives, then run
\code{make test}. The suite is cumulative and runs Chapter~1's and Chapter~2's tests too.

The rest will have worked solutions in the companion repository, published later as
\file{lectures/L03/appendix/e_solutions.md}, and the hand calculations have short answers in
Appendix~\ref{app:answers}. Work each one before you read it.

\textbf{No hardware is needed for any of this.}
\end{aside}

Each exercise is labelled with its kind. There is exactly one \textbf{Cross-check}, and here it is
Exercise~\ref{c3:ex:crosscheck}.

% ------------------------------------------------------------------------------------------
\exercise{The table}{Recall}
\label{c3:ex:table}
\begin{parts}
  \item How many vectors does the ATmega328P have, how many words does each slot occupy, and how
    many words of program memory does the whole table therefore take?
  \item RESET is in the table. What follows from that about the first instruction of your
    program?
  \item A slot is two words. Give the two things that can go in one, and say which this book uses.
  \item The datasheet says PCINT1 is at \code{0x0008}. What do you write after \code{.org}, and
    what is the name \code{m328Pdef.inc} gives that address? Assembly written for \code{avr-gcc}
    would need \code{.org 0x0010} for the same vector. Where does the factor of two come from?
  \item What is at word address \code{0x0007}, and why is it not a place anything may jump to?
\end{parts}

% ------------------------------------------------------------------------------------------
\exercise{Addresses and response times}{Hand calculation}
\label{c3:ex:response}
\begin{parts}
  \item Give the vector number, word address and byte address for PCINT0, PCINT2 and WDT.
  \item An interrupt flag is set at the exact moment an instruction finishes, and the vector holds
    an \code{rjmp}. How many cycles pass before the handler's first instruction? Show the terms.
  \item The same, but the flag is set just as a four-cycle \code{ret} begins.
  \item The compiler puts \code{jmp} in vector slots rather than \code{rjmp}. Redo (b) and (c) for
    a \code{jmp} table, and say in one sentence why a compiler would choose the slower one.
  \item Your handler measures 104 cycles including its \code{reti}, and the table uses
    \code{rjmp}. For how long are interrupts disabled? The answer is not 104.
\end{parts}

% ------------------------------------------------------------------------------------------
\exercise{The prologue}{Design}
\label{c3:ex:prologue}
A handler is about to be written. It calls \code{btn_pressed} and \code{led_toggle}, which
between them use \code{r18}, \code{r19}, \code{r24}, \code{X} and \code{Z}, and it uses
\code{r24} and \code{r25} itself to pass a pointer.
\begin{parts}
  \item List every register the handler must save, and say why the list is not ``the call-saved
    ones''.
  \item Write the prologue and epilogue, including \code{SREG}. Give the order precisely.
  \item Explain why the register is pushed \emph{before} \code{SREG} is read, and why \code{SREG}
    is restored \emph{before} the last \code{pop}.
  \item Somebody writes the epilogue as \code{pop r24}, \code{pop r25}, \code{out SREG, r24},
    \code{reti}. Say what is wrong with it, and describe the symptom in the interrupted program.
  \item How many cycles do your prologue and epilogue cost together? What fraction is that of a
    handler whose useful work is a single \code{sts}?
\end{parts}

% ------------------------------------------------------------------------------------------
\exercise{One vector, eight pins}{Design}
\label{c3:ex:onevector}
\begin{parts}
  \item Two buttons are wired to Arduino pins 12 and 13. How many interrupt vectors are involved,
    and which?
  \item One of them is pressed. Describe exactly what the handler is told.
  \item Write, in words, how the handler works out which button it was.
  \item The handler wants to know whether the button was \emph{pressed} or \emph{released}. Explain
    why it cannot find out from the pin alone, and what it would need to keep in order to know.
  \item A third button is added on Arduino pin 3. What changes in \code{PCICR}, what changes in the
    mask registers, and how many vectors are now involved?
\end{parts}

% ------------------------------------------------------------------------------------------
\exercise{The button driver}{Code}
\label{c3:ex:driver}
Write \file{drivers/source/btn.asm} to the specification in \secref{c3:sec:driver}.
\begin{parts}
  \item Write all six subroutines.
  \item Run \code{make test}.
  \item Deliberately swap the two return values in \code{btn_interrupt_enabled}, rebuild, and note
    which tests fail. One of them is not about \code{btn_interrupt_enabled} at all. Explain why it
    failed, then put the code back.
  \item Deliberately make \code{btn_disable_interrupt} clear \code{PCICR} as well as the mask bit.
    Say which test catches it, and why the test with two buttons does not.
\end{parts}

% ------------------------------------------------------------------------------------------
\exercise{The handler}{Code}
\label{c3:ex:handler}
Extend \file{drivers/app/main.asm} to the specification in \secref{c3:sec:handler}.
\begin{parts}
  \item Write the vector entry, the handler, the setup and the main loop. Put the LED on Arduino
    pin 13 and the button on pin 12.
  \item \code{make build} should build \file{app.hex}.
  \item Move \code{sei} from the end of setup to the beginning, before the structures are built.
    Say whether anything goes wrong in your program as it stands, and why. Then enable the
    button's interrupt before \code{btn_init} as well: describe what could now go wrong, and say
    why testing it would probably not show you.
  \item Put the button on pin 13 as well, so the LED and the button share a pin. Predict what
    happens before rebuilding, then check.
\end{parts}

% ------------------------------------------------------------------------------------------
\exercise{The shared counter}{Hand calculation}
\label{c3:ex:counter}
A handler increments a 16-bit counter in SRAM. The main loop reads it with two \code{lds}
instructions, low byte first.
\begin{parts}
  \item The counter holds \code{0x00FF} and the handler runs between the two \code{lds}. What value
    does the main loop assemble? Give it in hex and in decimal, and say what the counter actually
    held before and after.
  \item How many cycles wide is the window? At 16\,MHz, with the handler firing at 1\,kHz and the
    main loop reading once per millisecond, roughly how often would you expect a torn read?
  \item Write the fix, and give its cost in cycles.
  \item The obvious fix has a bug of its own. Say what it is and write the version without it.
  \item The same handler also sets an 8-bit flag that the main loop reads. Does that need
    protecting? Say why, in one sentence about instructions.
\end{parts}

% ------------------------------------------------------------------------------------------
\exercise[fillaccent]{How long is the door shut}{Cross-check}
\label{c3:ex:crosscheck}
\emph{Compute it by hand, measure it, reconcile.}

Your handler runs with interrupts disabled. This exercise is about finding out for how long, and
about which parts of that number you can measure and which you cannot.
\begin{parts}
  \item \textbf{By hand.} Count the cycles of your handler along the path where the button is
    \textbf{not} pressed: prologue, the call to \code{btn_pressed}, the compare, the branch taken,
    epilogue, \code{reti}. You measured \code{led_enabled} this way in Chapter~2, and
    \code{btn_pressed} measures the same way; if you have not measured it yet, do that now.
  \item \textbf{The whole window.} Add the entry to your handler's own cost: four cycles for the
    hardware's push, two for the \code{rjmp} in the vector slot, then the handler, then the
    \code{reti}. \textbf{This is the one number in the chapter you cannot check by measuring},
    because simavr charges nothing for the hardware's four-cycle push (\secref{c3:sec:sequence}).
    Write it down anyway, and expect (c) to disagree.
  \item \textbf{By measurement.}
\begin{shell}
make measure IMAGE=app \
    CALLS="led_init:0x0200:13 btn_init:0x0240:12 --drive B4=1 isr_pcint0"
\end{shell}
    \code{--drive B4=1} holds the button's pin high, which is a button that is \textbf{not}
    pressed. \code{IMAGE=app} because a handler lives in your program, not in the driver library.
    This works at all because \code{reti} pops a return address exactly as \code{ret} does.
  \item \textbf{Reconcile.} The two should agree. If they do not, the most likely culprits are the
    branch you charged one cycle instead of two, and forgetting that the measured figure already
    includes the \code{reti}.
  \item \textbf{Now the other path.} Predict the pressed-path cost before measuring it:
\begin{shell}
make measure IMAGE=app \
    CALLS="led_init:0x0200:13 btn_init:0x0240:12 --drive B4=0 isr_pcint0"
\end{shell}
    State the difference between the two paths and account for every cycle of it. There are five
    terms and two of them are negative.
  \item \textbf{The part you cannot measure.} Add the response time from \secref{c3:sec:sequence},
    not forgetting the one main-program instruction it guarantees after \code{reti}, to get the
    total worst-case delay before a \emph{second} interrupt can be served. Then set that
    against what the simulator does: it finishes the interrupted instruction as the device would,
    and then charges nothing at all for the push (\secref{c3:sec:sequence}). Note that there is no measurement
    to make here, because the harness offers none for this quantity, which is itself part of the
    answer. Say which of your sources you believe and why, and state the general rule you would
    apply next time a tool and a datasheet disagree.
  \item \textbf{What it means.} A second button on the same port is pressed while your handler is
    running. Is its interrupt lost or delayed? Now the \emph{same} button bounces twice during the
    window. Same question, and the answer is different. Explain both.
\end{parts}

% ------------------------------------------------------------------------------------------
\exercise{Why it responds to half your presses}{Design}
\label{c3:ex:bounce}
You build the program, wire a button, and press it. The LED toggles about half the time.
\begin{parts}
  \item Explain what is happening, in terms of the circuit rather than the code.
  \item Roughly how many interrupts does one press generate, and over what timescale?
  \item Why does ``about half'' make sense as an outcome, and what would the LED do if a press
    generated an odd number of interrupts instead?
  \item Give three ways to fix it, and say what each one costs. One of them needs something this
    book has not covered yet; name the chapter.
  \item Which of your three is wrong for a button and right for a rotary encoder? Say why.
\end{parts}
